\chapter{Project Design}
\label{ch:design}

This chapter turns the problem into requirements, models the system with context, data flow and decomposition diagrams, and explains the method together with the alternatives we rejected. It ends with the project plan and task allocation.

\section{Requirement Analysis}

\paragraph{Requirements elicitation.} Requirements were gathered from four sources: (i) the clinical guidelines that radiologists follow for nodules, namely the EUPS volume categories and the NELSON doubling-time categories~\cite{oudkerk2017eups,horeweg2014nelson}; (ii) meetings with our supervisor; (iii) the published limits of DTS and sparse-view reconstruction (\cref{ch:background}); and (iv) the constraints of a student project: public data only, no clinical scanner, limited computing resources and a fixed deadline. \Cref{tab:stakeholders} lists the stakeholders and their needs.

\begin{table}[!htbp]
\centering\small
\caption{Stakeholders and their needs.}
\label{tab:stakeholders}
\begin{tabularx}{\linewidth}{>{\raggedright\arraybackslash}p{2.8cm} >{\raggedright\arraybackslash}X >{\raggedright\arraybackslash}X}
\toprule
Stakeholder & Need & How MedLift-3D answers it \\
\midrule
Radiologist & A 3D size and growth rate from a low-dose exam, and to know when not to trust it & Mass, volume, category and doubling time, each with a calibrated interval and a low-confidence flag \\
Patient & Less radiation and cost per follow-up & DTS at about 0.1~mSv instead of CT \\
Low-resource hospital & Affordable equipment and software & DTS hardware; software runs on a single free or rented GPU \\
Examiner / researcher & Claims that can be checked & Every step ends in PASS/FAIL checks; test patients never used for tuning \\
\bottomrule
\end{tabularx}
\end{table}

\paragraph{Requirements analysis and negotiation.} Several requirements conflict; \cref{tab:conflicts} records how each conflict was resolved.

\begin{table}[!htbp]
\centering\small
\caption{Conflicting requirements and their resolution.}
\label{tab:conflicts}
\begin{tabularx}{\linewidth}{>{\raggedright\arraybackslash}p{3.4cm} >{\raggedright\arraybackslash}X >{\raggedright\arraybackslash}X}
\toprule
Conflict & Decision & Rationale \\
\midrule
Low dose vs accuracy & Simulate $I_0 = 2\times10^4$ photons per pixel and view (about 1/20--1/40 of clinical DTS dose) & Conservative: a method that works at this dose works at clinical dose \\
Clinical view count (60--74) vs compute & 15 views at equal total dose & Measured: the view count barely changes depth smear at equal dose (5.13 vs 5.34) \\
Fill the missing wedge vs never invent anatomy & Diffusion prior with X-ray data consistency, plus seen/unseen explanation & Every hypothesis agrees with the X-rays; what the prior adds is shown and tested \\
Single number vs honest uncertainty & Eight hypotheses + conformal calibration & Raw spread is not calibrated; conformal intervals have a coverage guarantee \\
Whole-volume accuracy vs GPU memory & Reconstruct only a 64~mm prism through the chest at the click & Keeps each ray's full path (a small box gave 73\% error) at $64{\times}128{\times}64$ voxels \\
Automatic detection vs false positives & The detector only suggests; the clinician clicks & The clinician stays in control; missed nodules can still be measured \\
\bottomrule
\end{tabularx}
\end{table}

\subsection{Functional and Nonfunctional Requirements}

Functional requirements state what the system must do, from loading the data to producing the report; nonfunctional requirements state the qualities it must have: speed, cost, reproducibility, validity, safety, usability, portability and privacy. Each requirement has a measurable acceptance criterion and the pipeline step that checks it (\cref{tab:fr,tab:nfr}). The criteria were chosen so that a step script can check them automatically, which makes the verification in \cref{ch:results} independent of anyone's judgement; \cref{tab:verify} reports the result for every row.

{\small
\begin{longtable}{>{\raggedright\arraybackslash}p{1.1cm} >{\raggedright\arraybackslash}p{5.6cm} >{\raggedright\arraybackslash}p{5.0cm} >{\raggedright\arraybackslash}p{1.6cm}}
\caption{Functional requirements.}\label{tab:fr}\\
\toprule
ID & Requirement & Acceptance criterion & Checked by \\
\midrule
\endfirsthead
\toprule
ID & Requirement & Acceptance criterion & Checked by \\
\midrule
\endhead
FR-01 & Load LIDC-IDRI CT and radiologist annotations; build a consensus mask ($\geq$50\% of readers); split by patient into train/cal/test & DICOM and pylidc slice positions agree; splits disjoint & Steps 02--04 \\
FR-02 & Simulate low-dose DTS: 15 views, $\pm15^\circ$, SID 1800~mm, 43~cm panel, Poisson noise & Detector integral equals object mass within 0.1\% in every view & Steps 01, 05 \\
FR-03 & Reconstruct the whole chest with SIRT, confined in depth to the patient & Residual below 0.9 of its start; missing wedge visible (depth/in-plane $>1.3$) & Step 06 \\
FR-04 & Suggest nodule candidates with a learned 3D detector & Test CPM higher than a training-free baseline & Step 14 \\
FR-05 & Reconstruct a 64~mm prism around a click & Prism keeps the column's attenuation (ratio $\approx 1$) & Step 07 \\
FR-06 & Generate eight X-ray-consistent CT hypotheses with the diffusion prior & Nodule-box PSNR of samples above the physics-only input on held-out pairs & Steps 08, 15 \\
FR-07 & Fit 3D Gaussians; measure mass, volume, 2D diameter and EUPS category & Finite values; failures reported, not hidden & Step 09 \\
FR-08 & Calibrate 90\% intervals and category probabilities & Test coverage $90\pm3\%$; Brier score below the raw ensemble vote & Steps 10, 16 \\
FR-09 & Measure growth on virtual follow-ups (VDT, MDT, NELSON categories) & Doubling-time error and category accuracy reported & Step 11 \\
FR-10 & Explain each result: seen/unseen split, detector saliency, error drivers & Tables and figures produced on test nodules & Step 17 \\
FR-11 & Interactive click-to-3D application and HTML report of all results & Demo self-test passes; report lists every check & Steps 12, 18 \\
\bottomrule
\end{longtable}
}

\begin{table}[!htbp]
\centering\small
\caption{Nonfunctional requirements.}
\label{tab:nfr}
\begin{tabularx}{\linewidth}{>{\raggedright\arraybackslash}p{3.2cm} >{\raggedright\arraybackslash}X >{\raggedright\arraybackslash}X}
\toprule
ID & Requirement & Acceptance criterion \\
\midrule
NFR-01 Performance & One click gives a result quickly & $\leq 60$~s per nodule on one GPU \\
NFR-02 Compute & The study and the application run on modest hardware & One GPU for the study; free GPU for the application \\
NFR-03 Reproducibility & Same input, same output; checkable steps & Fixed seed, one configuration file, PASS/FAIL checks, unit tests \\
NFR-04 Validity & No optimistic bias & Test patients unused for training or tuning; simulation and reconstruction on different grids \\
NFR-05 Safety & Never report a bare number & Every number has an interval; low-confidence flag when fewer than half of the hypotheses find the nodule or the spread exceeds $10\times$ \\
NFR-06 Usability & Usable by a clinician without coding & Browser interface, one click per nodule \\
NFR-07 Portability & Runs on common hardware & CPU smoke test; GPU with $\geq 24$~GB for training \\
NFR-08 Privacy & No identifiable patient data & Only de-identified public TCIA data \\
\bottomrule
\end{tabularx}
\end{table}

\subsection{Context Diagram}

The context diagram (\cref{fig:context}) depicts MedLift-3D as a single process with its external entities. The clinician is the only actor during run time: they send a case, a click and a depth and they receive the measurement report. The TCIA archive provides CT scans with the annotations from the radiologists and the cloud GPU provider performs the training and evaluation jobs and returns trained models and result, which are necessary only to build and evaluate the system, but not to use it.

\begin{figure}[htbp]
  \centering
  \resizebox{0.92\linewidth}{!}{% Context diagram (DFD level 0), Gane--Sarson notation
\begin{tikzpicture}
  \node[mainproc, minimum width=3.6cm, minimum height=2.0cm, text width=3.4cm] (sys) at (0,0)
       {\textbf{0}\\[2pt] \textbf{MedLift-3D}\\ click-to-3D nodule measurement};
  \node[entity] (clin) at (-7.3,0) {Clinician\\ (radiologist)};
  \node[entity] (tcia) at (0,3.4) {TCIA archive\\ (LIDC-IDRI CT + annotations)};
  \node[entity] (gpu)  at (0,-3.4) {Cloud GPU provider\\ (RunPod, Kaggle)};

  \draw[flow] ([yshift=4mm]clin.east) -- node[flowlbl, above=2pt] {case, click, depth} ([yshift=4mm]sys.west);
  \draw[flow] ([yshift=-4mm]sys.west) -- node[flowlbl, below=2pt] {mass, volume, 90\% intervals,\\ category probabilities, 3D view} ([yshift=-4mm]clin.east);
  \draw[flow] (tcia.south) -- node[flowlbl, right=2pt] {CT scans, 4-reader nodule annotations} (sys.north);
  \draw[flow] ([xshift=-8mm]sys.south) -- node[flowlbl, left=2pt] {training and evaluation jobs} ([xshift=-8mm]gpu.north);
  \draw[flow] ([xshift=8mm]gpu.north) -- node[flowlbl, right=2pt] {trained models, results} ([xshift=8mm]sys.south);
\end{tikzpicture}
}
  \caption{Context diagram (DFD level 0) in Gane--Sarson notation.}
  \label{fig:context}
\end{figure}

\subsection{Data Flow Diagram Level 1}

The level-1 DFD (\cref{fig:dfd1}) breaks down the system into eight processes and four data stores. The top row is run once for building the cases, training the models and calibrating, the second row is run every click and reads all four stores. The functional decomposition, entity-relationship model and the data dictionary are provided in \cref{sec:method}.

\begin{figure}[htbp]
  \centering
  \resizebox{\linewidth}{!}{% Data flow diagram level 1, Gane--Sarson notation
\begin{tikzpicture}[process/.append style={font=\scriptsize, text width=2.75cm, minimum width=2.95cm}, mainproc/.append style={font=\scriptsize, text width=2.75cm, minimum width=2.95cm}]
  \def\cA{0}\def\cB{4.2}\def\cC{8.4}\def\cD{12.6}\def\cE{16.8}
  \def\rTop{4.9}\def\rMid{2.45}\def\rBot{0}
  % offline row
  \node[entity, minimum width=2.6cm] (tcia) at (\cA,\rTop) {TCIA archive\\ (LIDC-IDRI)};
  \node[process] (p1) at (\cB,\rTop) {\textbf{1.0}\\ Build cases, consensus masks, split};
  \node[process] (p2) at (\cC,\rTop) {\textbf{2.0}\\ Simulate DTS and SIRT};
  \node[process] (p3) at (\cD,\rTop) {\textbf{3.0}\\ Train prior and detector};
  \node[process] (p4) at (\cE,\rTop) {\textbf{4.0}\\ Calibrate on cal.\ patients};
  % stores
  \datastore{d1}{(\cB,\rMid)}{D1}{Cases}
  \datastore{d2}{(\cC,\rMid)}{D2}{DTS + SIRT}
  \datastore{d3}{(\cD,\rMid)}{D3}{Models}
  \datastore{d4}{(\cE,\rMid)}{D4}{Calibration}
  % per-click row
  \node[entity, minimum width=2.6cm] (clin) at (\cA,\rBot) {Clinician};
  \node[process] (p5) at (\cB,\rBot) {\textbf{5.0}\\ Suggest candidates, select nodule};
  \node[process] (p6) at (\cC,\rBot) {\textbf{6.0}\\ Reconstruct prism};
  \node[mainproc] (p7) at (\cD,\rBot) {\textbf{7.0}\\ Sample prior, fit Gaussians};
  \node[process] (p8) at (\cE,\rBot) {\textbf{8.0}\\ Measure, calibrate, explain};

  \foreach \a/\b/\t in {tcia/p1/{CT}, p1/p2/{cases}, p2/p3/{pairs}, p3/p4/{models}} {
    \draw[flow] (\a) -- node[flowlbl, above=2pt] {\t} (\b);
  }
  \foreach \a/\b/\t in {clin/p5/{click}, p5/p6/{position}, p6/p7/{prism}, p7/p8/{samples}} {
    \draw[flow] (\a) -- node[flowlbl, below=2pt] {\t} (\b);
  }
  \foreach \p/\d/\t in {p1/d1/{cases}, p2/d2/{projections}, p3/d3/{weights}, p4/d4/{residuals}} {
    \draw[flow] (\p) -- node[flowlbl, right=1pt] {\t} (\d);
  }
  \foreach \d/\p/\t in {d1/p5/{case}, d2/p6/{DTS of case}, d3/p7/{prior weights}, d4/p8/{calibration}} {
    \draw[flow] (\d) -- node[flowlbl, right=1pt] {\t} (\p);
  }
  \draw[flow, rounded corners=5pt] (p8.south) -- ++(0,-0.65) -| node[flowlbl, pos=0.25, below] {report: mass and volume with 90\% intervals, category probabilities, seen/unseen split} (clin.south);
  \node[font=\scriptsize\itshape, text=quiet, anchor=west] at (-1.3,6.0) {offline, once};
  \node[font=\scriptsize\itshape, text=quiet, anchor=west] at (-1.3,1.1) {per click};
\end{tikzpicture}
}
  \caption{Data flow diagram level 1 (offline processes 1.0--4.0, per-click processes 5.0--8.0, data stores D1--D4).}
  \label{fig:dfd1}
\end{figure}





\subsection{UI Design}

The interface is a web-based application built on a single action, which is to click on the nodule (\cref{fig:ui}). The numbered parts are: (1) case and method selection; (2) the coronal DTS slice, with detector suggestions, annotated nodules (for evaluation) and with the click; (3) the depth slider, which provides the third coordinate; (4) the measurement report with mass and volume intervals, the EUPS category probabilities and the low-confidence flag; (5) a 3D view of the nodule and the prism input with its eight samples; and (6) a virtual follow-up with VDT and MDT.

\begin{figure}[htbp]
  \centering
  \resizebox{\linewidth}{!}{% User interface wireframe of the click-to-3D demo (Gradio)
\begin{tikzpicture}[font=\footnotesize,
    ctl/.style={draw=edge, fill=white, rounded corners=2pt, minimum height=0.55cm, anchor=west, inner xsep=4pt, font=\scriptsize},
    call/.style={circle, fill=accent, text=white, font=\scriptsize\bfseries, inner sep=1pt, minimum size=0.42cm}]
  % browser frame
  \draw[edge, fill=physics, rounded corners=4pt] (0,0) rectangle (16,10.2);
  \fill[edge!20, rounded corners=4pt] (0,9.6) rectangle (16,10.2);
  \foreach \x/\c in {0.3/red!60, 0.6/orange!70, 0.9/green!50!black} { \fill[\c] (\x,9.9) circle (0.08); }
  \node[anchor=west, font=\footnotesize\bfseries] at (1.3,9.9) {MedLift-3D: click-to-3D lung nodule assessment};
  % left panel
  \node[ctl] at (0.4,9.0) {Case: LIDC-IDRI-0018 $\blacktriangledown$};
  \node[ctl] at (3.9,9.0) {Method: diffusion + GS $\blacktriangledown$};
  \node[call] at (7.75,9.0) {1};
  \fill[black!75] (0.4,2.1) rectangle (7.4,8.4);
  \fill[black!15] (2.25,5.25) ellipse (1.25 and 2.4);
  \fill[black!15] (5.55,5.25) ellipse (1.25 and 2.4);
  \foreach \p in {(1.9,6.6),(5.0,4.0),(2.6,3.6),(6.0,6.9)} { \draw[orange!90!black, thick] \p circle (0.17); }
  \draw[green!60!black, thick] (5.7,6.2) circle (0.26);
  \draw[accent, very thick] (5.7,5.85) -- ++(0,0.7) (5.35,6.2) -- ++(0.7,0);
  \node[call] at (7.75,7.9) {2};
  \node[anchor=north west, font=\scriptsize, text=white] at (0.5,8.35) {coronal DTS slice};
  \node[anchor=south west, font=\scriptsize, text=white, align=left] at (0.5,2.15) {\textcolor{orange!80}{$\circ$} suggestion \quad \textcolor{green!70}{$\circ$} annotated \quad \textcolor{accentlight}{$+$} click};
  \draw[edge, thick] (0.6,1.4) -- (7.2,1.4);
  \fill[accent] (4.4,1.4) circle (0.13);
  \node[anchor=west, font=\scriptsize] at (0.5,0.85) {depth (slice) slider};
  \node[call] at (7.75,1.4) {3};
  % right panel: report
  \draw[edge, fill=white, rounded corners=3pt] (8.3,5.0) rectangle (15.6,9.3);
  \node[anchor=north west, font=\footnotesize\bfseries] at (8.45,9.2) {Measurement report};
  \node[anchor=north west, align=left, font=\scriptsize] at (8.45,8.7)
     {Mass $M$ [mg] with 90\% interval $[\ell, u]$\\[1pt]
      Volume $V$ [mm$^3$] with 90\% interval $[\ell, u]$\\[1pt]
      2D long-axis diameter [mm]\\[1pt]
      hypotheses that found the nodule: $k$ of 8};
  \node[anchor=west, font=\scriptsize\bfseries] at (8.45,6.55) {EUPS category probability};
  \foreach \i/\h/\t in {0/0.35/{$<$100}, 1/0.95/{100--300}, 2/0.45/{$\geq$300}} {
    \fill[accent!70] ({8.7+1.25*\i},5.35) rectangle ++(0.7,\h);
    \node[font=\scriptsize] at ({9.05+1.25*\i},5.15) {\t};
  }
  \node[draw=orange!90!black, text=orange!80!black, font=\scriptsize\bfseries, rounded corners=2pt] at (14.3,5.9) {low-confidence flag};
  \node[call] at (15.85,9.0) {4};
  % 3D view and prism strip
  \draw[edge, fill=white, rounded corners=3pt] (8.3,1.9) rectangle (11.8,4.7);
  \node[anchor=north west, font=\scriptsize\bfseries] at (8.4,4.65) {3D view};
  \shade[ball color=accent!60, opacity=0.8] (10.05,3.15) ellipse (0.9 and 0.6);
  \draw[edge, fill=white, rounded corners=3pt] (12.1,1.9) rectangle (15.6,4.7);
  \node[anchor=north west, font=\scriptsize\bfseries, text width=3.2cm] at (12.15,4.65) {prism: input + 8 samples};
  \foreach \i in {0,...,3} { \foreach \j in {0,1} { \fill[black!\the\numexpr 40+10*\i\relax] ({12.3+0.8*\i},{2.1+1.0*\j}) rectangle ++(0.7,0.85); } }
  \node[call] at (15.85,3.3) {5};
  % follow-up button
  \node[draw=accent, fill=accent, text=white, rounded corners=3pt, minimum width=7.0cm, minimum height=0.7cm, font=\scriptsize\bfseries] at (11.8,1.0) {Virtual follow-up (180 days): VDT, MDT};
  \node[call] at (15.85,1.0) {6};
\end{tikzpicture}
}
  \caption{User interface design of the click-to-3D application.}
  \label{fig:ui}
\end{figure}

\section{Detailed Methodology and Design}
\label{sec:method}

This section decomposes the system into its functions and data, derives what DTS can measure, and describes the proposed architecture, the network and the alternatives considered.

\subsection{Functional Decomposition Diagram}

The system decomposes into five subsystems (\cref{fig:fdd}). \emph{Data and simulation} selects the patients, downloads their CT, builds the cases and simulates DTS. \emph{Reconstruction} produces the whole-chest SIRT with depth support and the prism around a click. \emph{Learned models} contains the conditional diffusion prior with its data-consistency sampler, the nodule detector and the fit check. \emph{Measurement and UQ} fits Gaussians, measures mass, volume and diameter, and calibrates intervals and category probabilities. \emph{Interface and reporting} provides the click-to-3D application, the virtual follow-up, the explanations and the HTML report. Each function maps to one or more of the 18 scripted steps, so every box can be checked by its PASS/FAIL output.

\begin{figure}[htbp]
  \centering
  \resizebox{\linewidth}{!}{% Functional decomposition diagram
\begin{tikzpicture}[sub/.style={draw=edge, fill=accentlight, minimum width=3.3cm, minimum height=0.9cm, align=center, font=\footnotesize\bfseries, text width=3.1cm}]
  \node[mainproc, minimum width=5cm, text width=4.8cm, font=\small\bfseries] (root) at (7.8,2.0) {MedLift-3D system};
  \foreach \i/\n in {0/{1 Data and\\ simulation}, 1/{2 Reconstruction}, 2/{3 Learned models}, 3/{4 Measurement\\ and UQ}, 4/{5 Interface\\ and reporting}} {
    \node[sub] (s\i) at ({3.9*\i},0) {\n};
  }
  \draw[edge, thick] (root.south) -- ++(0,-0.6) coordinate (bus);
  \draw[edge, thick] (s0.north |- bus) -- (s4.north |- bus);
  \foreach \i in {0,...,4} { \draw[edge, thick] (s\i.north |- bus) -- (s\i.north); }
  \newcommand{\leaves}[5]{%
    \foreach \t [count=\k] in {{#2},{#3},{#4},{#5}} {
      \node[leaf] (l#1\k) at ($(s#1.south)+(0.3,{-0.95*\k+0.05})$) {\t};
      \draw[edge] ($(s#1.south west)+(0.12,0)$) |- (l#1\k.west);
    }}
  \leaves{0}{Select patients (LIDC)}{Download CT (TCIA)}{Cases + patient split}{Simulate 15-view DTS}
  \leaves{1}{Whole-chest SIRT}{Depth support}{Prism ROI-SIRT}{Background removal}
  \leaves{2}{Conditional diffusion prior}{SIRT data consistency}{Nodule detector}{Fit check}
  \leaves{3}{3D Gaussians}{Mass, volume, diameter}{Conformal intervals}{EUPS category probs.}
  \leaves{4}{Click-to-3D demo}{Virtual follow-up}{Seen/unseen, saliency}{HTML report}
\end{tikzpicture}
}
  \caption{Functional decomposition diagram: five subsystems and their functions.}
  \label{fig:fdd}
\end{figure}

\subsection{Entity-Relationship Model and Data Dictionary}

\Cref{fig:er} shows the data entities and their relationships: a patient has CT cases, a case contains nodules marked by three or four readers, and each nodule is measured by several methods, each measurement aggregating eight hypotheses that are fitted by Gaussian sets and calibrated per method. \Cref{tab:dict} gives the data dictionary.

\begin{figure}[htbp]
  \centering
  \resizebox{\linewidth}{!}{% Entity-relationship model (crow's-foot style cardinalities as labels)
\begin{tikzpicture}[font=\scriptsize,
    ent/.style={draw=edge, fill=white, rectangle split, rectangle split parts=2, rectangle split part fill={accentlight,white},
                align=left, text width=3.3cm, inner sep=3pt},
    rel/.style={edge, thick},
    card/.style={font=\scriptsize\bfseries, text=accent}]
  \node[ent] (pat) at (0,0) {\textbf{Patient}\nodepart{second}\underline{patient\_id}\\ split (train/cal/test)};
  \node[ent] (case) at (5.0,0) {\textbf{CT Case}\nodepart{second}\underline{case\_id}\\ patient\_id (FK)\\ HU volume\\ voxel geometry};
  \node[ent] (nod) at (10.0,0) {\textbf{Nodule}\nodepart{second}\underline{nid}\\ case\_id (FK)\\ consensus mask\\ diameter\\ GT volume, GT mass};
  \node[ent] (rdr) at (15.0,0) {\textbf{Reader Annotation}\nodepart{second}\underline{nid, reader}\\ reader mask\\ malignancy rating};
  \node[ent] (dts) at (5.0,-3.4) {\textbf{DTS Acquisition}\nodepart{second}\underline{case\_id}\\ 15 projections, angles\\ $I_0$, SIRT volume};
  \node[ent] (meas) at (10.0,-3.4) {\textbf{Measurement}\nodepart{second}\underline{nid, method}\\ volume, mass, diameter\\ log-spread $\hat\sigma$\\ interval $[\ell,u]$};
  \node[ent] (hyp) at (15.0,-3.4) {\textbf{Hypothesis}\nodepart{second}\underline{nid, method, k}\\ prism sample\\ volume, mass};
  \node[ent] (gau) at (15.0,-6.6) {\textbf{Gaussian Set}\nodepart{second}\underline{nid, method, k}\\ centres, scales,\\ rotations, amplitudes};
  \node[ent] (cal) at (10.0,-6.6) {\textbf{Calibration}\nodepart{second}\underline{method}\\ residuals, spreads\\ $\sigma$ floor};
  \node[ent] (mdl) at (5.0,-6.6) {\textbf{Model}\nodepart{second}\underline{name}\\ diffusion prior\\ or detector\\ weights, training log};

  \draw[rel] (pat) -- node[card, above, pos=0.15] {1} node[card, above, pos=0.85] {N} node[below, text=quiet] {has} (case);
  \draw[rel] (case) -- node[card, above, pos=0.15] {1} node[card, above, pos=0.85] {N} node[below, text=quiet] {contains} (nod);
  \draw[rel] (nod) -- node[card, above, pos=0.15] {1} node[card, above, pos=0.85] {3..4} node[below, text=quiet] {marked by} (rdr);
  \draw[rel] (case) -- node[card, right, pos=0.2] {1} node[card, right, pos=0.8] {1} node[left, text=quiet] {simulated as} (dts);
  \draw[rel] (nod) -- node[card, right, pos=0.2] {1} node[card, right, pos=0.8] {N} node[left, text=quiet] {measured by} (meas);
  \draw[rel] (meas) -- node[card, above, pos=0.15] {1} node[card, above, pos=0.85] {8} node[below, text=quiet] {aggregates} (hyp);
  \draw[rel] (hyp) -- node[card, right, pos=0.2] {1} node[card, right, pos=0.8] {1} node[left, text=quiet] {fitted by} (gau);
  \draw[rel] (cal) -- node[card, right, pos=0.2] {1} node[card, right, pos=0.8] {N} node[left, text=quiet] {calibrates} (meas);
  \draw[rel, rounded corners=4pt] (mdl.south) -- ++(0,-0.6) -| node[card, right, pos=0.97] {N} node[below, pos=0.25, text=quiet] {generates} ($(gau.east)+(0.45,0)$) |- (hyp.east);
  \node[card, right] at ($(mdl.south)+(0,-0.3)$) {1};
  \draw[rel] (dts) -- node[card, above, pos=0.15] {1} node[card, above, pos=0.85] {N} node[below, text=quiet] {input to} (meas);
\end{tikzpicture}
}
  \caption{Entity-relationship model of MedLift-3D (keys underlined, cardinalities in blue).}
  \label{fig:er}
\end{figure}

\begin{table}[!htbp]
\centering\small
\caption{Data dictionary.}
\label{tab:dict}
\begin{tabularx}{\linewidth}{>{\raggedright\arraybackslash}p{3.1cm} >{\raggedright\arraybackslash}p{3.0cm} >{\raggedright\arraybackslash}X >{\raggedright\arraybackslash}p{2.8cm}}
\toprule
Store / item & Format & Contents & Constraint \\
\midrule
D1 Case & int16 + packed bits + JSON header & CT in HU, consensus and four reader masks per nodule, voxel geometry, patient ID & HU in $[-1024, 3071]$ \\
Projections & float, $15{\times}536{\times}536$ & Noisy line integrals per view & Finite, max $<15$ \\
D2 SIRT volume & float, 2.5~mm voxels & Whole-chest DTS reconstruction & Patient depth + 5~mm \\
Prism & float, $64{\times}128{\times}64$, $1{\times}2{\times}1$~mm & Column through the chest at the click & Rays $\geq$50\% inside \\
Training pair & float16 & Condition (DTS prism) and target (CT prism) in HU & Train patients only \\
D3 Diffusion model & PyTorch checkpoint & 17.42~M parameters, EMA weights & Selected on held-out train pairs \\
D3 Detector & PyTorch checkpoint & 0.34~M parameters, heatmap 3D U-Net & Selected on cal.\ patients \\
Gaussians & float32 & Centres ($N{\times}3$), scales ($N{\times}3$), quaternions ($N{\times}4$), amplitudes ($N$) & $N \leq \num{30000}$ \\
Result row & JSON lines & Per nodule and method: volume, mass, log-spread, members, ground truth & One row per nodule \\
D4 Calibration & JSON & Per method: residuals, spreads, $\sigma$ floor & Cal.\ patients only \\
Report & JSON / UI & Mass and volume with intervals, category probabilities, diameter, low-confidence flag, run time & Interval type stated \\
\bottomrule
\end{tabularx}
\end{table}

\subsection{What DTS Can Measure}

Let $f(x,y,z)$ be the nodule's excess attenuation, with $y$ the depth axis, and let the view at tube angle $\theta$ project along a direction tilted by $\theta$ from $y$ in the $y$--$z$ plane, with detector coordinates $(x, v)$ and $v$ along $\mathbf{e}_\theta = (0, -\sin\theta, \cos\theta)$. The moments of the projection $p_\theta$ are
\begin{align}
  \iint p_\theta(x,v)\,\mathrm{d}x\,\mathrm{d}v &= M, \label{eq:m0}\\
  \iint v\, p_\theta(x,v)\,\mathrm{d}x\,\mathrm{d}v &= -\sin\theta\, M y_c + \cos\theta\, M z_c, \label{eq:m1}\\
  \iint (v-\bar v)^2 p_\theta\,\mathrm{d}x\,\mathrm{d}v &= \sin^2\theta\, S_{yy} - 2\sin\theta\cos\theta\, S_{yz} + \cos^2\theta\, S_{zz}, \label{eq:m2}
\end{align}
where $M$ is the mass, $(y_c,z_c)$ the centroid and $S$ the second central moments (for a cone beam, \cref{eq:m0} holds up to a known magnification factor). Every view gives the mass directly (\cref{eq:m0}). The depth position enters only through $\sin\theta \leq \sin 15^\circ$, making it $1/\sin15^\circ \approx 3.9$ times worse conditioned than in-plane position, and the depth extent $S_{yy}$, which volume needs, enters through $\sin^2\theta$, about $1/\sin^2 15^\circ \approx 15$ times worse. Step~01 confirms this numerically: the detector integral matched the mass to 0.03\% in all 15 views, and an object three times deeper with equal mass gave the same integrals.

\subsection{Proposed Architecture}

\Cref{fig:pipeline} gives an overview of the method; its stages are described below.

\begin{figure}[htbp]
  \centering
  \resizebox{\linewidth}{!}{% MedLift-3D method overview (research-style pipeline figure)
\begin{tikzpicture}[font=\footnotesize, >=Stealth,
    arr/.style={->, thick, edge},
    lbl/.style={align=center, font=\scriptsize, text width=3.1cm},
    stage/.style={font=\footnotesize\bfseries, text=accent},
    grp/.style={draw=edge!50, dashed, rounded corners=4pt, inner sep=6pt}]

  % cuboid: \cub{x}{y}{w}{h}{d}{fill}
  \newcommand{\cub}[6]{%
    \draw[edge, fill=#6] (#1,#2) rectangle ++(#3,#4);
    \draw[edge, fill=#6!70!white] (#1,#2+#4) -- ++(#5,0.55*#5) -- ++(#3,0) -- ++(-#5,-0.55*#5) -- cycle;
    \draw[edge, fill=#6!80!black] (#1+#3,#2) -- ++(#5,0.55*#5) -- ++(0,#4) -- ++(-#5,-0.55*#5) -- cycle;}

  % =================== row 1: acquisition and reconstruction (physics)
  \begin{scope}[on background layer]
    \fill[physics, rounded corners=4pt] (-0.3,4.35) rectangle (16.45,8.55);
  \end{scope}
  \node[anchor=north west, font=\small\bfseries, text=edge] at (-0.2,8.5) {A\quad Acquisition and reconstruction (physics)};

  % (1) projections: stack of views
  \foreach \i/\g in {3/55,2/45,1/35,0/25} {
    \draw[edge, fill=black!\g] ({0.25+0.16*\i},{5.15+0.16*\i}) rectangle ++(1.7,1.7);
  }
  \draw[white!80!black, thick] (0.6,6.0) ellipse (0.28 and 0.5);
  \draw[white!80!black, thick] (1.55,6.0) ellipse (0.28 and 0.5);
  \node[stage] at (1.35,7.75) {1};
  \node[lbl] at (1.35,4.75) {15 low-dose views\\ $\pm15^\circ$, $I_0{=}2{\times}10^4$};

  \draw[arr] (2.75,6.0) -- (3.35,6.0);

  % (2) whole-chest SIRT volume
  \cub{3.5}{5.15}{1.9}{1.7}{0.9}{white}
  \draw[accent, thick, fill=accent, fill opacity=0.15] (3.95,5.15) -- ++(0.9*0.55,0.55*0.9*0.55) -- ++(0,1.7) -- ++(-0.9*0.55,-0.55*0.9*0.55) -- cycle;
  \node[stage] at (4.75,7.75) {2};
  \node[lbl] at (4.75,4.75) {whole-chest SIRT\\ 2.5 mm voxels};

  \draw[arr] (6.45,6.0) -- (7.05,6.0);

  % (3) coronal slice + detector suggestions + click
  \draw[edge, fill=black!30] (7.2,5.15) rectangle ++(1.9,1.7);
  \draw[black!8, fill=black!8] (7.72,6.0) ellipse (0.36 and 0.62);
  \draw[black!8, fill=black!8] (8.58,6.0) ellipse (0.36 and 0.62);
  \foreach \p in {(7.6,6.35),(8.75,5.7),(7.85,5.55)} { \draw[orange!90!black, thick] \p circle (0.09); }
  \fill[accent] (8.5,6.3) circle (0.07);
  \draw[accent, thick] (8.5,6.05) -- ++(0,0.5) (8.25,6.3) -- ++(0.5,0);
  \node[stage] at (8.15,7.75) {3};
  \node[lbl] at (8.15,4.75) {detector suggests (\textcolor{orange!90!black}{$\circ$})\\ clinician clicks (\textcolor{accent}{$+$})};

  \draw[arr] (9.35,6.0) -- (9.95,6.0);

  % (4) prism through the chest
  \draw[edge] (10.15,5.15) rectangle ++(1.9,1.7);
  \draw[edge] (10.15,6.85) -- ++(1.2,0.66) -- ++(1.9,0) -- (12.05,6.85);
  \draw[edge] (13.25,7.51) -- ++(0,-1.7) -- (12.05,5.15);
  \draw[edge, densely dashed] (11.35,5.81) -- ++(1.9,0) (11.35,5.81) -- ++(0,1.7) (11.35,5.81) -- (10.15,5.15);
  \cub{10.85}{5.72}{0.5}{0.5}{1.2}{accentlight}
  \node[stage] at (11.45,7.75) {4};
  \node[lbl] at (11.45,4.75) {64 mm prism through\\ the chest ($64{\times}128{\times}64$)};

  % rows link
  \draw[arr, rounded corners=6pt] (13.5,6.0) -- (14.6,6.0) -- (14.6,4.05) -- (0.45,4.05) -- (0.45,2.05);
  \node[font=\scriptsize, fill=white, inner sep=1.5pt] at (8.0,4.05) {condition $\mathbf{c}$ = prism SIRT of the DTS};

  % =================== row 2: learned prior | measurement and uncertainty
  \begin{scope}[on background layer]
    \fill[learned, rounded corners=4pt] (-0.3,-0.75) rectangle (6.3,3.75);
    \fill[measure, rounded corners=4pt] (6.6,-0.75) rectangle (16.45,3.75);
  \end{scope}
  \node[anchor=north west, font=\small\bfseries, text=edge] at (0.8,3.7) {B\quad Learned prior};
  \node[anchor=north west, font=\small\bfseries, text=edge] at (6.7,3.7) {C\quad Measurement and calibrated uncertainty};

  % (5) conditional 3D U-Net
  \draw[edge, fill=accentlight] (1.25,0.75) -- (2.25,1.15) -- (2.25,2.05) -- (1.25,2.45) -- cycle;
  \draw[edge, fill=accent!45] (2.4,1.2) rectangle (2.9,2.0);
  \draw[edge, fill=accentlight] (3.05,1.15) -- (4.05,0.75) -- (4.05,2.45) -- (3.05,2.05) -- cycle;
  \draw[edge, dashed, ->] (1.75,2.35) to[bend left=30] (3.55,2.35);
  \node[font=\scriptsize, text=quiet] at (2.65,2.82) {skip connections};
  \node[align=center] at (0.45,1.6) {$\mathbf{x}_t$\\ $\oplus\,\mathbf{c}$};
  \draw[arr] (0.8,1.6) -- (1.2,1.6);
  \draw[arr] (4.1,1.6) -- (4.5,1.6);
  \node at (4.85,1.6) {$\hat{\mathbf{v}}_\theta$};
  \node[draw=accent, fill=white, rounded corners=2pt, font=\scriptsize\bfseries, text=accent] at (4.85,2.2) {$\times 8$};
  \draw[arr, accent, rounded corners=5pt] (4.85,1.3) -- (4.85,0.4) -- (0.45,0.4) -- (0.45,1.15);
  \node[font=\scriptsize, fill=learned, inner sep=1pt, text=accent!80!black] at (2.65,0.4) {DDIM step + SIRT data consistency};
  \node[stage] at (5.6,3.0) {5};
  \node[lbl, text width=5cm] at (2.65,-0.25) {conditional 3D diffusion U-Net\\ (17.4\,M parameters)};

  \draw[arr] (5.2,1.6) -- (6.95,1.6);

  % (6) 3D Gaussians
  \draw[edge, fill=white] (7.0,0.6) rectangle ++(2.0,2.0);
  \foreach \cx/\cy/\rx/\ry/\r in {7.8/1.65/0.42/0.22/30, 8.15/1.35/0.35/0.18/-20, 8.0/1.85/0.25/0.15/70, 8.4/1.75/0.22/0.12/10} {
    \fill[accent, opacity=0.35, rotate around={\r:(\cx,\cy)}] (\cx,\cy) ellipse (\rx cm and \ry cm);
  }
  \node[stage] at (8.0,3.0) {6};
  \node[lbl, text width=2.4cm] at (8.0,0.0) {3D Gaussians\\ closed-form mass};

  \draw[arr] (9.05,1.6) -- (9.45,1.6);

  % (7) measurements
  \node[draw=edge, fill=white, rounded corners=3pt, align=left, text width=1.75cm, minimum height=2.0cm, font=\scriptsize] at (10.45,1.6)
    {mass $M$\\[2pt] volume $V$\\[2pt] diameter $d$\\[2pt] spread $\sigma$};
  \node[stage] at (10.45,3.0) {7};
  \node[lbl, text width=2.2cm] at (10.45,0.0) {8 measurements\\ per nodule};

  \draw[arr] (11.5,1.6) -- (11.9,1.6);

  % (8) conformal predictive distribution
  \begin{scope}[shift={(12.05,0.6)}]
    \draw[->, edge] (0,0) -- (2.1,0) node[right, font=\scriptsize] {$V$};
    \draw[->, edge] (0,0) -- (0,2.0) node[left, font=\scriptsize] {$F$};
    \fill[accent, opacity=0.15] (0.55,0) rectangle (1.45,1.85);
    \draw[accent, very thick, domain=0:1.95, samples=60, smooth] plot (\x, {1.8/(1+exp(-6*(\x-1.0)))});
    \draw[edge, densely dashed] (0.55,0) -- (0.55,1.85) (1.45,0) -- (1.45,1.85);
    \node[font=\scriptsize] at (0.55,-0.2) {$\ell$};
    \node[font=\scriptsize] at (1.45,-0.2) {$u$};
  \end{scope}
  \node[stage] at (13.1,3.0) {8};
  \node[lbl, text width=3.2cm] at (13.1,0.0) {conformal predictive\\ distribution};

  \draw[arr] (14.4,1.6) -- (14.65,1.6);

  % (9) report
  \draw[edge, fill=white] (14.7,0.5) -- ++(1.65,0) -- ++(0,1.9) -- ++(-0.3,0.3) -- ++(-1.35,0) -- cycle;
  \node[align=left, font=\scriptsize] at (15.5,1.5) {$M\ [\ell,u]$\\ $V\ [\ell,u]$\\ $P(\mathrm{cat.})$\\ seen/unseen};
  \node[stage] at (15.5,3.0) {9};
  \node[lbl, text width=1.8cm] at (15.5,0.0) {report};
\end{tikzpicture}
}
  \caption[MedLift-3D method overview]{MedLift-3D method overview. (A) Physics: low-dose DTS is simulated, reconstructed and cut to a prism around the click. (B) A conditional 3D diffusion prior generates eight hypotheses of the prism, each kept consistent with the X-rays. (C) Gaussians give closed-form mass; the eight measurements are calibrated into intervals and category probabilities.}
  \label{fig:pipeline}
\end{figure}

\paragraph{Stage 1: simulated low-dose DTS.} The CT table is removed, then 15 views over $\pm15^\circ$ are simulated at CT resolution with 0.5~mm ray sampling, SID 1800~mm, a 43~cm flat panel with 0.8~mm pixels, monochromatic attenuation and Poisson noise at $I_0 = 2\times10^4$. The geometry matches the GE VolumeRAD sweep; reconstruction uses different grids than simulation to avoid the inverse crime.

\paragraph{Stage 2: whole-chest SIRT with depth support.} SIRT (40 iterations, 2.5~mm voxels) is confined in depth to the patient (body outline + 5~mm). Without this support the reconstruction placed attenuation in the air in front of and behind the patient, and on real chests the prism lost about a third of its attenuation; with it the prism keeps 98.7\%.

\paragraph{Stage 3: learned detector.} A small 3D U-Net~\cite{cicek2016unet} (0.34~M parameters) predicts a nodule heatmap from the SIRT volume, and its peaks are offered as suggestions. A Laplacian-of-Gaussian and Hessian blob detector is the training-free baseline.

\paragraph{Stage 4: prism ROI.} Around the click, a $64\times64$~mm column through the whole chest ($64\times128\times64$ voxels at $1\times2\times1$~mm) is reconstructed by SIRT using only rays that carry at least 50\% of their attenuation inside the prism; the rest of the chest is subtracted using the coarse reconstruction. Depth is kept whole because each ray's attenuation is spread along its full path: a small 3D box around the nodule gave a 73\% error against 5--7\% for the prism.

\paragraph{Stage 5: conditional 3D diffusion prior.} A 3D U-Net (\cref{fig:unet}; base 32 channels, multipliers 1/2/4/4, two residual blocks per level, attention at the lowest level, 17.42~M parameters) learns $p(\text{CT prism}\mid\text{DTS prism})$ with v-prediction, a cosine schedule over 1000 steps, EMA weights and bf16 mixed precision. Sampling uses 50 DDIM steps; from 60\% of the way through, every denoised estimate $\hat{\mathbf{x}}_0$ receives two SIRT data-consistency iterations,
\begin{equation}
  \hat{\mathbf{x}}_0 \leftarrow \hat{\mathbf{x}}_0 + \lambda\, C A^{\top} R\,(\mathbf{b} - A\hat{\mathbf{x}}_0),
\end{equation}
before the next DDIM step. Eight samples give eight hypotheses that agree with the X-rays but differ where the X-rays are blind.



\paragraph{Stage 6: 3D Gaussians and measurement.} For each hypothesis the nodule box ($40\times50\times40$~mm) is fitted with anisotropic Gaussians (100 fitting and 100 refinement iterations against the measured projections, the rest of the prism fixed). A seeded half-maximum segmentation, confined to a plausible nodule extent, gives the volume; \cref{eq:gaussmass} gives the mass; and the long-axis 2D diameter is kept as the clinical baseline.

\paragraph{Stage 7: uncertainty.} The eight hypotheses give a log-scale spread $\hat\sigma$. On the calibration patients this spread is calibrated with (a) split-conformal asymmetric intervals (step~10) and (b) normalised and Mondrian conformal predictive systems with crepes (step~16), which give full predictive distributions, EUPS category probabilities and category sets.

\paragraph{Stage 8: growth.} Each nodule is grown in the CT by a known factor, its projections are updated and fresh noise added; VDT and MDT are measured and assigned to the NELSON categories.

\paragraph{Stage 9: explanation.} The true CT prism is split into seen and unseen parts by conjugate-gradient least squares (CGLS)~\cite{hestenes1952cg} on the prism projection operator. We measure how much of the nodule's mass and shape lies in each part and the prior's \emph{unseen skill}, the correlation of its unseen content with the truth's (\cref{fig:seen}). Input$\times$gradient saliency explains the detector, and rank correlations relate the measurement error to nodule size, depth and contrast.

\subsection{Diffusion Network Architecture}

\Cref{fig:unet} shows the conditional 3D U-Net. The noisy prism and the DTS condition enter as two channels; four resolution levels with 32, 64, 128 and 128 channels use two residual blocks each in the encoder and three in the decoder, with skip connections between them; multi-head self-attention runs at the coarsest level ($8{\times}16{\times}8$) and in the bottleneck; and a sinusoidal timestep embedding passed through an MLP is added inside every residual block. The output layer is initialised to zero so that training starts from a stable state.

\begin{figure}[htbp]
  \centering
  \resizebox{0.95\linewidth}{!}{% Conditional 3D U-Net of the diffusion prior (medlift3d/diffusion/unet3d.py)
\begin{tikzpicture}[font=\footnotesize, >=Stealth,
    res/.style={draw=edge, fill=accentlight, minimum width=#1, inner sep=0pt},
    att/.style={draw=accent, fill=accent!45, minimum width=0.22cm, inner sep=0pt},
    down/.style={->, thick, red!65!black},
    up/.style={->, thick, green!45!black},
    skip/.style={->, densely dashed, edge},
    lab/.style={font=\scriptsize, align=center}]

  % level geometry: y centre, block height (resolution), block width (channels)
  \def\yA{6.0}\def\yB{4.4}\def\yC{2.8}\def\yD{1.2}
  \def\xe{3.2}\def\xd{10.0}

  % ----- encoder
  \node[res=0.45cm, minimum height=1.3cm] (e0) at (\xe,\yA) {};
  \node[res=0.65cm, minimum height=1.0cm] (e1) at (\xe,\yB) {};
  \node[res=0.85cm, minimum height=0.75cm] (e2) at (\xe,\yC) {};
  \node[res=0.85cm, minimum height=0.55cm] (e3) at (\xe,\yD) {};
  \node[att, minimum height=0.55cm, right=0pt of e3] (a3) {};
  \node[lab, left=4pt of e0] {$2\times$ ResBlock, 32 ch\\ $64{\times}128{\times}64$};
  \node[lab, left=4pt of e1] {$2\times$ ResBlock, 64 ch\\ $32{\times}64{\times}32$};
  \node[lab, left=4pt of e2] {$2\times$ ResBlock, 128 ch\\ $16{\times}32{\times}16$};
  \node[lab, left=4pt of e3] {$2\times$ (ResBlock + attention)\\ 128 ch, $8{\times}16{\times}8$};
  \draw[down] (e0.south) -- (e1.north);
  \draw[down] (e1.south) -- (e2.north);
  \draw[down] (e2.south) -- (e3.north);

  % ----- bottleneck
  \node[res=0.55cm, minimum height=0.55cm] (m1) at (6.25,-0.2) {};
  \node[att, minimum height=0.55cm, right=1pt of m1] (ma) {};
  \node[res=0.55cm, minimum height=0.55cm, right=1pt of ma] (m2) {};
  \node[lab, below=3pt of ma] {bottleneck: ResBlock -- attention -- ResBlock};
  \draw[->, thick, edge] (a3.south) |- (m1.west);

  % ----- decoder
  \node[res=0.85cm, minimum height=0.55cm] (d3) at (\xd,\yD) {};
  \node[att, minimum height=0.55cm, right=0pt of d3] (b3) {};
  \node[res=0.85cm, minimum height=0.75cm] (d2) at (\xd,\yC) {};
  \node[res=0.65cm, minimum height=1.0cm] (d1) at (\xd,\yB) {};
  \node[res=0.45cm, minimum height=1.3cm] (d0) at (\xd,\yA) {};
  \draw[->, thick, edge] (m2.east) -| (d3.south);
  \draw[up] (d3.north) -- (d2.south);
  \draw[up] (d2.north) -- (d1.south);
  \draw[up] (d1.north) -- (d0.south);
  \node[lab, right=4pt of b3] {$3\times$ (ResBlock + attention)};
  \node[lab, right=4pt of d2] {$3\times$ ResBlock, 128 ch};
  \node[lab, right=4pt of d1] {$3\times$ ResBlock, 64 ch};
  \node[lab, right=4pt of d0] {$3\times$ ResBlock, 32 ch};

  % ----- skips
  \foreach \e/\d in {e0/d0, e1/d1, e2/d2} { \draw[skip] (\e.east) -- node[above, lab, text=quiet]{concatenate} (\d.west); }
  \draw[skip] (a3.east) -- node[above, lab, text=quiet]{concatenate} (d3.west);

  % ----- input / output
  \node[draw=edge, fill=white, rounded corners=2pt, lab] (in) at (\xe,8.0) {$\mathbf{x}_t \oplus \mathbf{c}$\\ 2 channels};
  \draw[->, thick, edge] (in.south) -- node[right, lab]{conv $3^3$} (e0.north);
  \node[draw=edge, fill=white, rounded corners=2pt, lab] (out) at (\xd,8.0) {$\hat{\mathbf{v}}_\theta$\\ 1 channel};
  \draw[->, thick, edge] (d0.north) -- node[left, lab]{GN, SiLU, conv $3^3$} (out.south);

  % ----- time embedding
  \node[draw=edge, fill=white, rounded corners=2pt, lab] (t) at (6.6,8.0) {timestep $t$ $\rightarrow$ sinusoidal + MLP\\ \textcolor{accent}{$\mathbf{e}_t$ added in every ResBlock}};

  % ----- legend
  \begin{scope}[shift={(0.6,-1.9)}]
    \node[res=0.35cm, minimum height=0.3cm] at (0,0) {}; \node[lab, anchor=west] at (0.3,0) {ResBlock (GN, SiLU, conv, $+\,\mathbf{e}_t$)};
    \node[att, minimum height=0.3cm] at (5.4,0) {}; \node[lab, anchor=west] at (5.6,0) {multi-head self-attention};
    \draw[down] (9.4,0.15) -- (9.4,-0.15); \node[lab, anchor=west] at (9.55,0) {stride-2 conv};
    \draw[up] (11.8,-0.15) -- (11.8,0.15); \node[lab, anchor=west] at (11.95,0) {upsample + conv};
  \end{scope}
\end{tikzpicture}
}
  \caption[Conditional 3D U-Net of the diffusion prior]{Conditional 3D U-Net of the diffusion prior. The noisy prism $\mathbf{x}_t$ is concatenated with the DTS condition $\mathbf{c}$; the timestep embedding is added in every residual block; the network predicts the velocity $\hat{\mathbf{v}}_\theta$.}
  \label{fig:unet}
\end{figure}

\subsection{Alternative Solutions}

 \Cref{tab:alternatives} lists the alternatives and why they were rejected. The initial design of this project generated whole $512^3$ volumes with a latent diffusion model; it is included in the table because its lessons shaped the final design: a large generator judged by training losses could not show whether a single nodule measurement was right, so the goal was narrowed to clinical numbers that can be verified.

\begin{table}[!htbp]
\centering\small
\caption{Alternative solutions and the rationale for the selected design.}
\label{tab:alternatives}
\begin{tabularx}{\linewidth}{>{\raggedright\arraybackslash}p{3.4cm} >{\raggedright\arraybackslash}X >{\raggedright\arraybackslash}p{3.6cm}}
\toprule
Alternative & Why rejected & Selected instead \\
\midrule
FBP or SIRT alone & Cannot fill the missing wedge; physics-only volume errors of 80--90\% on small nodules & SIRT as input and data-consistency step, plus a learned prior \\
Neural radiance fields~\cite{mildenhall2020nerf} & Slow per-patient optimisation, implicit geometry, no prior for unseen depth & Explicit 3D Gaussians with closed-form mass \\
Whole-volume latent diffusion at $512^3$ (initial design) & Too large to train and verify with our resources; evaluated by generative losses, not clinical numbers & Small conditional prior on a prism, judged by calibrated measurements \\
Small 3D box ROI & Each ray's attenuation lies mostly outside the box: 73\% error vs 5--7\% & Prism through the whole chest \\
GAN prior & Training instability; one output per input & Diffusion: stable training, many hypotheses \\
Unconditional prior with posterior sampling only & Slower and weaker for a strongly ill-posed \ang{30} problem & Conditional model + SIRT data consistency \\
Raw ensemble variance or MC dropout & Not calibrated; no coverage guarantee & Ensemble spread calibrated by conformal prediction \\
Volume-only reporting & Volume depends on the poorly measured depth & Mass first, volume second, both with intervals \\
\bottomrule
\end{tabularx}
\end{table}

\section{Project Plan}

MedLift-3D was developed over twelve months in an iterative research and development cycle, organised in four phases that match the task groups of the Gantt chart in \cref{fig:gantt}. The iterative approach suits the experimental nature of generative AI: the prior, the measurement and the uncertainty layer were revised whenever an intermediate verification step showed a problem.

\paragraph{Phase 1: Inception, data and feasibility (months 1--3).} \emph{Objective:} establish the theoretical basis and a reliable data pipeline. \emph{Tasks:} 1.1 literature review and ethics, 1.2 dataset acquisition and preprocessing, 1.3 DTS simulation, 1.4 baseline reconstruction. \emph{Deliverables:} literature review, curated LIDC-IDRI cases with consensus masks, the DTS simulator and baseline SIRT results.

\paragraph{Phase 2: Core generative architecture (months 3--6).} \emph{Objective:} build and train the learned prior. \emph{Tasks:} 2.1 projection physics and its verification (including the projection-moment analysis), 2.2 prism ROI reconstruction, 2.3 conditional diffusion prior with data consistency, 2.4 diffusion training. \emph{Deliverables:} verified differentiable projector, prism reconstruction, trained diffusion prior.

\paragraph{Phase 3: Gaussians, detection and uncertainty (months 6--9).} \emph{Objective:} turn reconstructions into calibrated, explainable measurements. \emph{Tasks:} 3.1 Gaussian parameterisation and closed-form mass, 3.2 nodule detector, 3.3 conformal uncertainty, 3.4 explainability. \emph{Deliverables:} measurement module, detector, calibrated intervals and probabilities, seen/unseen analysis.

\paragraph{Phase 4: Validation, interface and deployment (months 9--12).} \emph{Objective:} evaluate the complete system and deliver it. \emph{Tasks:} 4.1 pilot and full-scale evaluation, 4.2 growth and sensitivity studies, 4.3 user interface development, 4.4 documentation and defense. \emph{Deliverables:} results of the 583-patient study, the click-to-3D application, the final report.

\begin{sidewaysfigure}
  \centering
  \resizebox{0.95\textheight}{!}{% FYDP-III schedule (weeks of the trimester)
\begin{ganttchart}[
    hgrid, vgrid={*{8}{dotted}},
    x unit=1.15cm, y unit chart=0.52cm,
    title height=1, title label font=\footnotesize\bfseries,
    bar/.append style={fill=accentlight, draw=accent},
    bar label font=\footnotesize, bar height=0.55,
    group/.append style={fill=edge!70, draw=edge}, group label font=\footnotesize\bfseries,
    group peaks height=.2, group peaks tip position=0, group left shift=0, group right shift=0, group peaks width=0.2,
    milestone/.append style={fill=accent, draw=accent}, milestone left shift=-0.5, milestone right shift=-0.5, milestone label font=\footnotesize\itshape,
    link/.style={->, thick, edge}]{1}{9}
  \gantttitle{FYDP-III trimester (weeks)}{9} \\
  \gantttitlelist{1,...,9}{1} \\
  \ganttgroup{Analysis and design}{1}{3} \\
  \ganttbar{Problem re-analysis, projection moments}{1}{2} \\
  \ganttbar{Requirements, design, verification plan}{2}{3} \\
  \ganttgroup{Implementation}{2}{8} \\
  \ganttbar[name=phys]{Physics, DTS simulation, unit tests}{2}{4} \\
  \ganttbar[name=rec]{SIRT, prism ROI, 3D Gaussians}{3}{5} \\
  \ganttbar[name=dif]{Diffusion prior and training}{4}{6} \\
  \ganttbar{Detector, conformal UQ, growth}{5}{7} \\
  \ganttbar{Explainability, demo, HTML report}{6}{8} \\
  \ganttgroup{Verification and study}{6}{9} \\
  \ganttbar[name=pilot]{Real-data pilots (4 and 40 patients)}{6}{7} \\
  \ganttmilestone[name=full]{Full 583-patient run}{8} \\
  \ganttbar{Report, slides, defense}{8}{9} \\
  \ganttmilestone{Final submission}{9}
  \ganttlink{phys}{rec}
  \ganttlink{rec}{dif}
  \ganttlink{pilot}{full}
\end{ganttchart}
}
  \caption{Gantt chart of the twelve-month project schedule, by phase and task.}
  \label{fig:gantt}
\end{sidewaysfigure}

\section{Task Allocation}

Responsibilities were allocated according to the technical area in which each member is strongest, so that the critical paths of theoretical modelling, system optimisation and application engineering could progress in parallel. The mathematically intensive parts (the diffusion prior, its conditioning and data consistency, and the statistical calibration) need a solid theoretical grounding and were led by the lead AI architect. The systems engineer focused on performance and infrastructure, making the volumetric reconstruction feasible in GPU memory and running the computation on cloud GPUs. The full-stack developer connected the algorithms to clinical use, ensuring data quality, the user interface and the evaluation outputs. Every pipeline step has one owner who is accountable for its PASS/FAIL checks, and all members reviewed every step before the full run. \Cref{tab:tasks} gives the detailed allocation.

\begin{table}[!htbp]
\centering\small
\caption{Detailed task allocation matrix.}
\label{tab:tasks}
\begin{tabularx}{\linewidth}{>{\raggedright\arraybackslash}p{3.2cm} >{\raggedright\arraybackslash}p{3.0cm} >{\raggedright\arraybackslash}X}
\toprule
Team member & Role and focus & Detailed responsibilities \\
\midrule
Jubair Ahmed (011222136) & Lead AI architect (theoretical modelling) &
\textbullet~Projection-moment analysis of what DTS can measure \newline \textbullet~Design and mathematical formulation of the conditional diffusion prior (v-prediction, cosine schedule, DDIM sampling) \newline \textbullet~SIRT data-consistency sampling to suppress anatomical hallucination \newline \textbullet~Conformal intervals and calibrated category probabilities \\
\midrule
Foyez Ahammed Naeem (011222011) & Systems engineer (optimisation and backend) &
\textbullet~Differentiable DTS projector and memory-efficient prism reconstruction \newline \textbullet~GPU optimisation of training and sampling (mixed precision, batch sizing) \newline \textbullet~3D Gaussian fitting and closed-form mass; nodule detector training \newline \textbullet~Cloud infrastructure (RunPod CPU and GPU pods, Kaggle), the full run and backups \\
\midrule
Md Motasim Billah (011222016) & Full-stack developer (data and visualisation) &
\textbullet~Data curation and preprocessing of LIDC-IDRI (selection, consensus masks, DTS simulation inputs) \newline \textbullet~Clinician-facing click-to-3D interface with interactive 3D visualisation \newline \textbullet~Automated evaluation pipelines: accuracy, growth and explainability outputs, HTML report \\
\midrule
All members & -- & Literature review, verification of every PASS/FAIL step, report writing and presentation \\
\bottomrule
\end{tabularx}
\end{table}

\section{Summary}

The requirements were elicited from guidelines, the supervisor and published limits, negotiated where they conflicted, and specified with measurable acceptance criteria. The design measures mass first, reconstructs a prism around the click, fills depth with a data-consistent diffusion prior, measures with Gaussians, and calibrates and explains every number. \Cref{ch:results} verifies each requirement.
